\section{Introduction}
\label{sec:intro}

Removing transformer blocks and then continuing pretraining (CPT, often called
``healing'') is an attractive route to smaller language models
~\citep{gromov2024unreasonable,shortenedllama,men2024shortgpt,
yang2024laco}. Held-out perplexity is a natural training signal, but using its
improvement as a certificate for a target capability makes a stronger
measurement assumption: that likelihood and the target evaluation recover
together after the structural intervention.

Prior compression work already shows that aggregate language-model metrics can
miss changes on knowledge-sensitive, downstream, or safety evaluations
~\citep{costcompression,jaiswal2024truth,beyondperplexity}, while depth-pruning
studies report recovery curves and loss--task gaps
~\citep{gromov2024unreasonable,shortenedllama,minitron,iterabre}. We do not claim
either phenomenon as new. We instead study proxy validity in one densely
measured OLMo setting: \emph{is in-domain perplexity sufficient to certify
performance on the measured knowledge-sensitive evaluations along the literal
observed recovery paths?}

We instantiate the study on OLMo-2-1124-7B~\citep{olmo2}. The principal
construction retains the first 14 of 32 pretrained blocks, appends two fresh
blocks, and continues pretraining for 200k optimizer steps. Operating points
include the intact base; full32 continued pretraining through an available 25k
checkpoint; a frozen-prefix construction; a fully random same-shape model; and
a non-contiguous ShortGPT-16 construction. Evaluation combines same-source
held-out perplexity, answer-letter and complete-option MMLU, and no-retrieval
PopQA, TriviaQA, and NQ-open.

\begin{figure*}[t]
\centering
\includegraphics[width=\textwidth]{figures/fig1_pruneheal_autofigure.pdf}
\caption{\textbf{Proxy-validity view of the observed paths.}
\textbf{Left:} the single keep14 run improves on both measured axes from 128k
to 200k steps, while answer-letter MMLU remains far below the intact base.
\textbf{Right:} endpoint and null operating points make three boundaries
visible: random initialization has an answer-letter chance score but a high
content-score floor; ShortGPT reaches a stronger endpoint under a coupled
construction; and full32 ends at 25k, not 200k. PPL is evaluated on a
same-source held-out Dolmino shard. Each trained construction is one run;
token counts are nominal presentations. The figure supports proxy
insufficiency on these literal paths, not a causal, localization, deletion, or
universal-dynamics claim.}
\label{fig:proxy-validity}
\end{figure*}

Figure~\ref{fig:proxy-validity} summarizes the measurement result. In the
principal run, in-domain PPL improves through 200k, but answer-letter MMLU and
all three closed-book scores remain well below the base. Complete-option MMLU
raises the apparent score, yet the random-init content baseline shows that this
interface has a substantial non-target floor. ShortGPT is stronger on both PPL
and MMLU, so the keep14 endpoint is construction-dependent rather than a
property of nominal 16-layer depth. Because ShortGPT changes four structural
factors together, and random/frozen change learning rate or trainable modules,
these comparisons bound alternative explanations without isolating causes.

\paragraph{Protocol boundaries at first use.}
Every trained construction is a single run. full32 is observed only through
25k, whereas the principal and same-shape endpoints use 200k. Prefix depths
have unequal retained checkpoints selected after target metrics were inspected;
there was no registered common stopping rule, and equal optimizer steps are not
equal FLOPs across depths. ShortGPT couples selection, inherited count,
final-block retention, and fresh-tail use. Random-init and frozen-front are
operating points, not clean ablations. All reported PPL is in-domain on a
disjoint shard from the same Dolmino/DCLM source.

\paragraph{Contributions.}
\begin{itemize}[leftmargin=1.2em,itemsep=2pt]
\item \textbf{A literal observed-path test of proxy sufficiency.} We jointly
track in-domain likelihood and knowledge-sensitive evaluations without
interpreting one training realization as a general recovery law.
\item \textbf{Operating points that bound alternatives.} A short-horizon intact
branch, same-shape null points, two MMLU interfaces, closed-book generation, and
a coupled ShortGPT construction constrain corpus-shift, interface, and
construction-only accounts while keeping confounds visible.
\item \textbf{A reporting recommendation, not a new pruning method.} The case
motivates separate reporting of likelihood, target evaluation, interface,
exact construction, budget/compute, and run-level uncertainty.
\end{itemize}
